\section{Implementation overview}
\label{sec:glsp-implementation-overview}

This section provides a structured overview of the UVL GLSP implementation, covering both the server and the client components.
The server is described first, as it owns the domain model, the GLSP session state, and the operation handlers that implement editing behavior.
The client is then presented as the VS~Code activation layer responsible for launching the server and connecting the custom editor to it.
The presented implementation details are reduced to the essentials that demonstrate how the GLSP runtime is extended to support UVL-specific editing logic.
Some representative code snippets are included to illustrate key points, but the full implementation is available in the accompanying code repository for readers who want to explore further~\cite{TODO_sourcecode}.

\subsection{Server bootstrap}

The server entry point is \texttt{UVLServerLauncher}, which parses command-line arguments, instantiates the diagram module, and starts the socket-based GLSP server.
The launcher is intentionally minimal as can be seen in Listing~\ref{lst:server-launcher-overview}, delegating all configuration logic to the module system.
All concrete behavior is assembled through the module system.

\lstinputlisting[
    language=Java,
    float,
    floatplacement=ht,
    caption={[Server launcher and module bootstrap]Server launcher and module bootstrap.},
    label={lst:server-launcher-overview}
]{code/glsp_overview/UVLServerLauncher.java}

\texttt{UVLDiagramModule} serves as the central configuration point of the server.
It binds the diagram configuration, model state, source model storage, graphical model factory, layout engine, and all operation handlers into the GLSP dependency injection framework.
In effect, this module is where UVL-specific editing logic is plugged into the generic GLSP runtime.
In Listing~\ref{lst:diagram-module-overview}, there are examples of bindings for the source model storage, the model state, and the operation handlers.
The full module contains additional bindings for other components such as the graphical model factory and layout engine~\cite{TODO_sourcecode}.

\lstinputlisting[
    language=Java,
    float,
    floatplacement=ht,
    caption={[Essential bindings in the diagram module]Essential bindings in the diagram module.},
    label={lst:diagram-module-overview}
]{code/glsp_overview/UVLDiagramModule.java}

Each registered handler translates an incoming GLSP operation into a targeted mutation of both the UVL domain model and its graphical representation.
Handlers that correspond to generic editing actions, such as label editing and deletion, subclass existing GLSP framework classes and augment them to additionally update the UVL feature model.
For example, \texttt{UVLApplyLabelEditOperationHandler} extends the framework's \texttt{GModelApplyLabelEditOperationHandler}, adding a step that propagates the name change into the domain model.
Handlers for UVL-specific operations, such as creating features, attributes, and relation edges, implement the handler interface directly and encode entirely UVL-specific creation logic with no framework base class to build on.

\texttt{UVLModelState} is the interface that aggregates the active feature model, the GModel root, and the model index into a single injectable dependency. 
\texttt{UVLModelStateImpl} is its bound implementation, registered in \texttt{UVLDiagramModule} and injected into every handler that requires access to the model.
Handlers never receive the index directly, they always obtain it through \texttt{UVLModelState\#getIndex()}, keeping the index lifecycle coupled to the session state.

\subsection{Source model and indexing}

Persistence is managed by \texttt{UVLSourceModelStorage}, which maintains a paired representation of each diagram. 
On load, it reads the textual \texttt{.uvl} file into an in-memory feature model and simultaneously reads the adjacent notation file, which shares the same base name but carries the extension \texttt{.notation.json}.
The notation file encodes the GLSP graph state, including element identifiers and layout coordinates, as a JSON document.

As discussed in Section~\ref{sec:concept_glsp}, a UVL file should contain only data relevant to the feature tree or the behavioral program.
Layout data, such as element positioning, could technically be encoded as feature attributes, but conflating notation with semantics would invalidate the file for any non-GLSP tooling that expects a semantically clean UVL document.
Therefore, the paired-file design keeps the two concerns strictly separate, ensuring that the \texttt{.uvl} file remains consumable by the broader UVL ecosystem regardless of whether the graphical editor was involved.
Both files are written together in a single save step so that the user perceives a single coherent artifact.

\lstinputlisting[
    language=Java,
    float,
    floatplacement=ht,
    caption={[Source model loading and paired notation file]Source model loading and paired notation file.},
    label={lst:source-model-overview}
]{code/glsp_overview/UVLSourceModelStorage.java}

Because GLSP graph elements are identified by opaque string IDs while the UVL domain model is a structured object graph, handlers cannot navigate from a selected diagram element to the corresponding feature or constraint without an explicit mapping.
\texttt{UVLModelIndex} provides this mapping by maintaining a bidirectional association between element IDs and UVL domain objects.
The index is rebuilt whenever the model root is replaced and cleared when the client session is disposed, preventing stale references.

\lstinputlisting[
    language=Java,
    float,
    floatplacement=ht,
    caption={[Rebuilding and querying the model index]Rebuilding and querying the model index.},
    label={lst:model-index-overview}
]{code/glsp_overview/UVLModelIndex.java}

\subsection{Extensibility example: \texttt{BPDiagramModule}}

The GLSP module system makes it straightforward to provide an alternative diagram configuration for behavioral programming variant by supplying a custom diagram module.
In a separate package, a \texttt{BPDiagramModule} can simply subclass the existing UVL diagram module and override the same binding points used above.
The Listing \ref{lst:bp-diagram-module-overview} shows how the BP variant replaces the \texttt{SourceModelStorage} and the \texttt{UVLModelState} implementation while reusing the common bindings from \texttt{UVLDiagramModule}.

\lstinputlisting[
    language=Java,
    float,
    floatplacement=ht,
    caption={[Example \texttt{BPDiagramModule} overriding core bindings]Example \texttt{BPDiagramModule} overriding core bindings.},
    label={lst:bp-diagram-module-overview}
]{code/glsp_overview/BPDiagramModule.java}

In this pattern, \texttt{BPSourceModelStorage} applies the parsing logic using the BP grammar variant, while \texttt{BPModelStateImpl} can extend \texttt{UVLModelStateImpl} to use the \texttt{BPFeatureModel} metamodel class.
Addionally, the operation handlers are overridden to implement BP-specific editing logic, as well as adding new handlers for BP-specific operations such as creating b-threads and b-events.
Because all bindings are resolved through the module, the replacement is localised to the module class and requires no invasive changes to the GLSP framework or to existing handlers.

\subsection{Client integration via VS~Code}

As described in Section~\ref{sec:glsp}, the standard GLSP approach separates client rendering functionality from IDE-specific integration code, so that the rendering package can be reused across different IDE frontends.
This separation is applied here as well.
The VS~Code integration package depends on \texttt{uvl-sprotty} as a shared rendering library, and its sole responsibility is to establish the connection between the VS~Code webview runtime and that shared package.
The main entry point of the VS~Code integration is \texttt{UvlDiagramStarter}, which initializes the shared \texttt{UvlDiagramContainer} with all relevant rendering configuration.

\lstinputlisting[
  language=Java,
  float,
  floatplacement=ht,
  caption={[VS~Code webview integration and container setup]VS~Code webview integration and container setup.},
  label={lst:client-integration-overview}
]{code/glsp_overview/UvlDiagramStarter.ts}

The \texttt{resolveWebviewPluginModules} function collects all additional container modules declared under \texttt{containerModules} in the \texttt{profiles.json} configuration file.
These plugin modules extend the base container with supplementary packages from \texttt{uvl-sprotty}.
Command contributions, by contrast, are injected at a separate point in the activation code.

\subsection{Client rendering package: \texttt{uvl-sprotty}}

The rendering logic shared across all UVL GLSP client variants is contained in \texttt{client/packages/uvl-sprotty}.
Its central file, \texttt{uvl-diagram-module.ts}, is the single location where all UVL-specific UI behavior is bound into the runtime.

\lstinputlisting[
  language=Java,
  float,
  floatplacement=ht,
  caption={[Excerpt from \texttt{uvl-diagram-module.ts} showing the registered UVL UI elements]Excerpt from \texttt{uvl-diagram-module.ts} showing the registered UVL UI elements.},
  label={lst:uvl-diagram-module-overview}
]{code/glsp_overview/uvl-diagram-module.ts}

Comparing this listing with the server-side \texttt{UVLDiagramModule} shown earlier reveals a structural symmetry that is intentional.
The GLSP framework applies an inversion of control pattern based on dependency injection on both sides of the client--server boundary.
On the server, bindings are managed by Google Guice~\cite{noauthor_googleguice_2026}.
On the client, the same pattern is realised through InversifyJS~\cite{noauthor_inversifymonorepo_2026}.
In both cases, every service and component is registered in a global DI container and can be straightforwardly extended or replaced with a custom implementation, making the architecture uniformly composable across the full stack.

This composability is expressed concretely through the element registrations in the module.
Each call to \texttt{configureModelElement} associates a UVL type constant with a model class and a view class.
When the container receives a model element of type \texttt{UVLModelTypes.FEATURE}, for instance, it instantiates a \texttt{FeatureNode} as the model object and delegates rendering to \texttt{FeatureNodeView}, which produces the corresponding SVG output.
The same pattern applies uniformly to attributes, relation edges, and other model elements, keeping the registration declarative and the mapping between type, model, and view explicit.

Beyond element registration, the module also wires interaction and layout behavior into the container.
For example, the \texttt{CenteredAnchor} overrides the default anchor computation to center edge connections between parent and child features.
The \texttt{HighlightElementsActionHandler} receives highlight actions dispatched from the server and applies them to the graphical view.
Although the handler is implemented generically and could be reused by other server components, in this implementation it is primarily used to highlight b-events selected in the behavioral program view.

Overall, the implementation tries to apply the principle of separation of concerns consistently at every level of the architecture.
At the framework level, the shared DI container model means that both the server and the client are uniformly composable.
Individual services can be replaced or extended without restructuring the surrounding system.
At the component level, all modeling logic is concentrated on the server, leaving the client as a thin wiring layer that could be substituted with any other GLSP-compatible frontend without affecting the server.
At the file level, layout data is kept out of the \texttt{.uvl} file, ensuring that the source model remains valid for the broader UVL ecosystem independently of the graphical editor.